%==============================================================================
\section{Discussion}
\label{sec:discussion}
%==============================================================================

\subsection{Relationship to Inference-Time Graph Methods}
\label{sec:discuss-inference}

A growing body of work leverages code structure graphs at inference time to improve LLM-based software engineering.
CodexGraph~\cite{liu2025codexgraph} materializes a code property graph in a Neo4j database that LLM agents query via Cypher; RepoGraph~\cite{ouyang2025repograph} provides line-level code graph navigation as a plug-in for SE agents; LocAgent~\cite{chen2025locagent} and RANGER~\cite{shah2025ranger} employ graph-guided multi-hop reasoning and MCTS-based traversal for fault localization and retrieval.
GraphCoder~\cite{liu2024graphcoder} and ARISE~\cite{seddik2026arise} incorporate data-flow edges but remain limited to intra-procedural scope.

We highlight two axes along which our work occupies a distinct position.
First, \textbf{when structural information enters the pipeline}: inference-time methods construct and query graphs per request, incurring latency proportional to graph complexity.
Our approach amortizes all graph computation to training time (CodeQL extraction costs $\sim$8.9 seconds per instance, $\sim$110 CPU-hours for 44K instances), after which the trained encoder operates at standard embedding-lookup speed with zero graph overhead.
Second, \textbf{depth of data-flow analysis}: most inference-time methods use structural relations limited to call graphs and import dependencies, whereas we leverage CodeQL's bounded inter-procedural taint tracking (2-hop call-graph scope), capturing cross-function data-flow paths.

The two paradigms are complementary: a model trained with structure-grounded negatives can serve as a first-stage retriever whose embeddings already encode data-flow semantics, reducing the search space for inference-time graph agents that perform targeted reranking.

\subsection{When Does Mining Strategy Matter?}
\label{sec:discuss-when}

Our results suggest that mining strategy effectiveness is not universal but depends on encoder properties and evaluation conditions.

\paragraph{Architecture-scale boundary condition.}
UniXcoder-base (125M, encoder-only) exhibits negligible strategy differentiation: all four strategies produce R@1 within \placeholder{0.4pp} of each other.
This contrasts with the 3B decoder-only encoders, where strategy effects are measurable.
We interpret this as a capacity threshold: smaller encoders may lack the representational flexibility to exploit the additional signal in harder negatives.
The boundary between ``mining-sensitive'' and ``mining-insensitive'' encoders likely depends on both parameter count and pretraining quality, but our data (\placeholder{$n{=}3$} encoders) is insufficient to characterize the threshold precisely.

\paragraph{Representation geometry hypothesis.}
Among the 3B encoders, Qwen2.5-Coder-3B (positive-pair cosine similarity 0.986) and StarCoder2-3B (0.798) respond differently to mining strategies.
Negative difficulty margin analysis shows that Qwen's near-saturated embedding space leaves minimal room for discrimination (random margin = 0.074), while StarCoder maintains 2.6$\times$ larger margins.
We offer this as a candidate heuristic---encoders with pos\_sim $>$ 0.98 may benefit most from the hardest available negatives---while noting that $n{=}2$ encoders at the same scale provides suggestive evidence, not a general law.

\subsection{Evaluation Methodology Implications}
\label{sec:discuss-eval}

Our finding that checkpoint selection criterion yields a \placeholder{+16.3pp} R@1 difference---exceeding all strategy effects---has practical implications for the community.

\paragraph{Reporting recommendations.}
We recommend that contrastive retrieval papers explicitly report: (1) the checkpoint selection criterion used, (2) results under both best-by-loss and best-by-metric checkpoints as a minimum transparency standard, and (3) evaluation at multiple pool sizes to characterize the generalization of observed effects.
The dual-checkpoint protocol we adopt (reporting $\theta^*_{\text{R@1}}$ as primary with $\theta^*_{\text{loss}}$ as robustness check) provides a concrete template.

\paragraph{Reinterpreting prior findings.}
The magnitude of the criterion effect raises questions about prior work that reports negative mining results without specifying checkpoint selection.
If prior studies selected checkpoints by validation loss---a reasonable default from the language modeling tradition---they may have systematically underestimated the benefit of harder negatives, which tend to produce loss--R@1 divergence more frequently (as observed in our training dynamics analysis, \S\ref{sec:training-dynamics}).

\subsection{Limitations}
\label{sec:limitations}

\paragraph{Encoder diversity.}
We evaluate three encoders spanning two architectures (encoder-only, decoder-only) and two scales (125M, 3B).
This is sufficient for initial observations about architecture-scale boundaries, but insufficient for characterizing general thresholds.
Encoder-only models at 3B scale and decoder-only models at 125M remain untested, confounding architecture with scale.

\paragraph{Single benchmark.}
All training and primary evaluation uses SWE-smith-derived data from Python repositories.
Cross-benchmark validation (e.g., RepoBench-R) is planned but not yet complete.
The generalization of our findings to other languages, domains, or retrieval benchmarks is uncertain.

\paragraph{Static mining.}
All negative mining is performed once before training and remains fixed throughout.
Dynamic re-mining~\cite{xiong2021ance}---where negatives are refreshed using the current encoder's embeddings---could yield progressively harder negatives as the model improves.
We do not evaluate this variant and acknowledge the static-vs-dynamic confound as a limitation.

\paragraph{Structural extraction coverage.}
CodeQL extraction achieves \placeholder{90\%} structural data rate post-reextraction; the remaining \placeholder{10\%} falls back to beniget (AST-level) or embedding-distance negatives.
While high, this means a fraction of dataflow-aware training instances use degraded structural signal.

\paragraph{No gradient-based comparison.}
We compare static mining strategies but do not implement gradient-based methods (e.g., CoCoHaNeRe's differentiable mechanism).
Our claim is that static embedding mining provides a reliable baseline, not that it matches or exceeds gradient-based approaches.

\subsection{Future Work}
\label{sec:future}

\paragraph{Hybrid retrieval pipelines.}
Integrating a structure-trained encoder as the first-stage retriever within inference-time graph systems could combine embedding efficiency with runtime graph precision.

\paragraph{Representation-guided strategy selection.}
Our pos\_sim threshold observation, if validated across more encoders, could enable automated strategy selection without exhaustive ablation: compute representation geometry statistics on a calibration set, then select the appropriate mining strategy.

\paragraph{Dynamic and multi-lingual extension.}
Adapting the pipeline to dynamically re-mine negatives during training, and extending CodeQL queries beyond Python, would address two current limitations.

%==============================================================================
